{\bf [40 points] Dimensionality Reduction}\\

In this question, we will study and compare two widely used dimensionality reduction methods
in single-cell RNA-seq analysis: (i) Principal Components Analysis (PCA) and
(ii) t-distributed Stochastic Neighbor Embedding (t-SNE).
The goal is to understand how the objective functions of these methods determine whether
they preserve global or local structure in high-dimensional data.

\vspace{10pt}

Throughout this problem, assume that the data matrix
$X \in \mathbb{R}^{n \times d}$ has been mean-centered
(i.e., each feature has zero mean).


\vspace{10pt}
\textbf{Your task}
\vspace{10pt}



% {\bf Principal Components Analysis}\\

% In this part, we study Principal Components Analysis (PCA) as a linear dimensionality
% reduction method. 
Assume that the data matrix
$X \in \mathbb{R}^{n \times d}$ is \textbf{mean-centered}.

\begin{enumerate}

% =================================================
\item (5 points) \textbf{PCA: Variance maximization formulation.}

Consider projecting each data point $x_i \in \mathbb{R}^d$ onto a one-dimensional
direction $w \in \mathbb{R}^d$ with $\|w\|_2 = 1$, producing
\[
z_i = w^\top x_i.
\]

Show that the variance of the projected data can be written as
\[
\mathrm{Var}(z) = w^\top \Sigma w,
\]
where $\Sigma = \frac{1}{n} X^\top X$ is the sample covariance matrix.

%%%%%%%%%%%%%%%%%%
\begin{solution}

\end{solution}
%%%%%%%%%%%%%%%%%%

% =================================================
\item (5 points) \textbf{PCA: Eigenvalue problem.}

PCA seeks the direction $w$ that maximizes $w^\top \Sigma w$ subject to $\|w\|_2 = 1$.

\begin{enumerate}
    \item Using a Lagrange multiplier, show that the optimal solution satisfies
    \[
    \Sigma w = \lambda w.
    \]
    \item Explain why the eigenvector corresponding to the largest eigenvalue
    is selected as the first principal component.
\end{enumerate}


%%%%%%%%%%%%%%%%%%
\begin{solution}

\end{solution}
%%%%%%%%%%%%%%%%%%

% =================================================
\item (5 points) \textbf{PCA and global structure.}

PCA can also be viewed as minimizing the total squared reconstruction error
\[
\sum_{i=1}^{n} \|x_i - \hat{x}_i\|_2^2,
\]
where $\hat{x}_i$ denotes the projection of $x_i$ onto the low-dimensional
subspace.

Explain why this objective makes PCA particularly sensitive to \emph{global}
structure in the data. In your answer, discuss how distances between far-apart
points influence the solution.

%%%%%%%%%%%%%%%%%%
\begin{solution}

\end{solution}
%%%%%%%%%%%%%%%%%%
% {\bf [t-SNE: 20 points] t-distributed Stochastic Neighbor Embedding}\\

% In this part, we study t-distributed Stochastic Neighbor Embedding (t-SNE), a nonlinear
% dimensionality reduction method designed to preserve local neighborhood structure.

\vspace{15pt}

% \begin{enumerate}

% =================================================
\item (8 points) \textbf{ t-SNE: High-dimensional similarities}

In t-SNE, similarities between points in the original space are defined as conditional
probabilities:
\[
p_{j|i} =
\frac{\exp\left(-\|x_i - x_j\|^2 / (2\sigma_i^2)\right)}
{\sum_{k \neq i} \exp\left(-\|x_i - x_k\|^2 / (2\sigma_i^2)\right)}.
\]

Explain the intuition behind this definition. What does it mean for $p_{j|i}$ to be
large or small?

%%%%%%%%%%%%%%%%%%
\begin{solution}

\end{solution}
%%%%%%%%%%%%%%%%%%

% =================================================
% \item (7 points) \textbf{t-SNE: Low-dimensional similarities and heavy-tailed distribution}

% In the low-dimensional embedding $\{y_i\}_{i=1}^n$, t-SNE defines similarities as
% \[
% q_{ij} =
% \frac{(1 + \|y_i - y_j\|^2)^{-1}}
% {\sum_{k \neq l} (1 + \|y_k - y_l\|^2)^{-1}}.
% \]

% Explain why a heavy-tailed distribution (such as the Student-$t$ distribution)
% is used for defining similarities in the low-dimensional space.

% %%%%%%%%%%%%%%%%%%
% \begin{solution}

% A heavy-tailed distribution assigns non-negligible probability to moderately large
% distances in low-dimensional space. This helps alleviate the crowding problem,
% where many points that are moderately far apart in high dimension are forced too
% close together in two or three dimensions. Using a Student-$t$ distribution allows
% dissimilar points to be placed farther apart without their similarity values becoming
% too small, leading to better separation of clusters in the visualization.

% \end{solution}
% %%%%%%%%%%%%%%%%%%

% % =================================================
% \item (8 points) \textbf{t-SNE: KL divergence and local structure in}

% t-SNE finds the low-dimensional embedding by minimizing the Kullback--Leibler (KL)
% divergence between the high-dimensional and low-dimensional similarity distributions:
% \[
% \mathcal{L} = \sum_{i,j} p_{ij} \log \frac{p_{ij}}{q_{ij}}.
% \]

% Explain why this objective emphasizes preserving \emph{local} neighborhood structure.
% In your explanation, comment on the asymmetry of the KL divergence and the role of
% pairs with large $p_{ij}$.

% %%%%%%%%%%%%%%%%%%
% \begin{solution}

% The KL divergence $\mathrm{KL}(P\|Q)$ is asymmetric and weights each term by $p_{ij}$.
% As a result, pairs of points that are close neighbors in the high-dimensional space
% have large $p_{ij}$ values and dominate the loss. If such pairs are mapped far apart
% in the low-dimensional space (so that $q_{ij}$ is small), the penalty becomes large.
% In contrast, pairs that are far apart in high dimension typically have
% $p_{ij} \approx 0$ and contribute negligibly to the loss regardless of their
% low-dimensional distance. Thus, the objective focuses on preserving local
% neighborhood relationships rather than global geometry.

% \end{solution}
%%%%%%%%%%%%%%%%%%

% \end{enumerate}
% =================================================
\item (12 points) \textbf{t-SNE: KL divergence, asymmetry, and locality}

In t-SNE, similarities in the low-dimensional embedding are defined as
\[
q_{ij} =
\frac{(1 + \|y_i - y_j\|^2)^{-1}}
{\sum_{k \neq l} (1 + \|y_k - y_l\|^2)^{-1}}.
\]

The low-dimensional embedding is obtained by minimizing the Kullback--Leibler (KL)
divergence between the high-dimensional and low-dimensional similarity distributions:
\[
\mathcal{L} = \sum_{i,j} p_{ij} \log \frac{p_{ij}}{q_{ij}}.
\]

\begin{enumerate}
    \item Explain how the weighting by $p_{ij}$ affects which pairs of points
    contribute most to the loss.
    \item  Consider a pair of points that are far apart in the high-dimensional
    space but close together in the low-dimensional embedding. Explain why this
    typically does \emph{not} incur a large penalty.
    \item Qualitatively describe what would change if t-SNE minimized
    $\mathrm{KL}(Q\|P)$ instead of $\mathrm{KL}(P\|Q)$.
\end{enumerate}

%%%%%%%%%%%%%%%%%%
\begin{solution}


\end{solution}
%%%%%%%%%%%%%%%%%%

\vspace{20pt}

% =================================================
% {\bf (10 points) PCA vs. t-SNE}\\
\item (5 points) \textbf{Comparison: PCA vs. t-SNE}

\begin{quote}
Why is PCA generally considered a \emph{global} dimensionality reduction method,
while t-SNE is considered a \emph{local} dimensionality reduction method?
\end{quote}

Your answer should explicitly reference the objective functions of both methods.

%%%%%%%%%%%%%%%%%%
\begin{solution}

\end{solution}
%%%%%%%%%%%%%%%%%%

\end{enumerate}
